\subsection{Disk--finite-set calipers}

Let $\mathsf J$ denote rotation through $\pi/2$.  For a centered disk
$\mathcal D_\eta$ and a finite point set $P=\{p_1,\ldots,p_m\}$, put
\[
 \Psi_P(n)=
 \sum_{j=0}^2
 \max\left\{\eta,\max_i\langle p_i,\mathsf R^j n\rangle\right\}.
\]

\begin{proposition}[Disk--finite-set caliper theorem]
\label{prop:new-disk-finite-caliper}
If an orientation exists in which all point projections in the three support
directions are at most $\eta$, then
\[
 \Lambda(\mathcal D_\eta\cup P)=\frac{3\eta}{h}.
\]
Otherwise a minimizing normal may be chosen from the finite list
\[
 \pm\frac{\mathsf J(p_k-p_i)}{\|p_k-p_i\|}
 \quad(p_i\ne p_k)
\]
and
\[
 \frac{\eta p_i\pm
 \sqrt{\|p_i\|^2-\eta^2}\,\mathsf Jp_i}{\|p_i\|^2}
 \quad(\|p_i\|\ge\eta).
\]
Equivalently, one side of a least enclosing triangle either contains two
points, or is tangent to the disk and contains one point.
\end{proposition}

\begin{proof}
The tie directions partition the normal circle into open cells with fixed
support sources.  If $q$ supports come from the disk, then on one cell
\[
 \Psi_P(n(\theta))=q\eta+\langle v,n(\theta)\rangle
\]
for a fixed vector $v$.  If any point is active, then
$\langle v,n(\theta)\rangle>0$, and hence
\[
 \Psi_P''(n(\theta))=-\langle v,n(\theta)\rangle<0.
\]
Thus no such cell has an interior minimum.  A boundary is a point--point or
point--disk tie.  On a disk-only cell the value is $3\eta$, the global
minimum.
\end{proof}

\subsection{Universal terminal G: complementary gap}

\begin{lemma}[Disk-plus-point formula]
\label{lem:new-disk-point-formula}
Let $\mathcal D_\eta$ be the disk of radius $\eta$ about $O$, and let
$\|G\|=\rho$.  Then
\[
 \Lambda(\mathcal D_\eta\cup\{G\})=
 \begin{cases}
  3\eta/h,&\rho\le2\eta,\\[2mm]
  \dfrac{3\eta+\sqrt{3(\rho^2-\eta^2)}}{2h},&\rho>2\eta.
 \end{cases}
 \tag{6.27}
\]
\end{lemma}

\begin{figure}[!htbp]
\centering
\begin{minipage}[t]{.485\linewidth}
\centering
\resizebox{.96\linewidth}{!}{%
\begin{tikzpicture}[x=3cm,y=3cm,>=Latex]
  \node[panellabel,anchor=west] at (-.72,.93) {(a) $\rho\le2\eta$};
  \coordinate (O) at (0,0);
  \coordinate (G) at (.55,0);
  \coordinate (Z) at (.70,0);
  \draw[CenterBlue!35,fill=CenterBlue!8,line width=.8pt] (O) circle (.35);
  \draw[black,line width=1pt]
    (Z)--(-.35,.6062)--(-.35,-.6062)--cycle;
  \fill (O) circle (.012) node[below left=1pt,figlabel] {$O$};
  \fill[DemandRed] (G) circle (.018) node[above=2pt,figlabel] {$G$};
  \fill (Z) circle (.012) node[right=2pt,figlabel] {$Z$};
  \draw[dimline] (0,-.13)--(.35,-.13)
    node[midway,below=1pt,figlabel] {$\eta$};
  \draw[dimline] (0,.13)--(.70,.13)
    node[midway,above=1pt,figlabel] {$2\eta$};
  \node[figlabel,align=center] at (0,-.86)
    {$G\in[O,Z]$ and the disk-only minimizer still contains $G$\\
     $\displaystyle\Lambda=\frac{3\eta}{h}$};
\end{tikzpicture}%
}
\end{minipage}\hfill
\begin{minipage}[t]{.485\linewidth}
\centering
\resizebox{.96\linewidth}{!}{%
\begin{tikzpicture}[x=2.55cm,y=2.55cm,>=Latex]
  \node[panellabel,anchor=west] at (-.78,1.05) {(b) $\rho>2\eta$};
  \coordinate (O) at (0,0);
  \coordinate (G) at (1.2,0);
  \coordinate (T) at (.1021,.3348);
  \coordinate (A) at (-.4778,.5116);
  \coordinate (B) at (-.0819,-1.1972);
  \draw[CenterBlue!35,fill=CenterBlue!8,line width=.8pt] (O) circle (.35);
  \draw[black,line width=1pt] (A)--(B)--(G)--cycle;
  \draw[DemandRed,line width=1pt] (O)--(G);
  \draw[RadialTeal,line width=.9pt] (O)--(T);
  \draw[GapPurple,line width=.9pt] (T)--(G);
  \fill (O) circle (.012) node[below left=1pt,figlabel] {$O$};
  \fill[DemandRed] (G) circle (.018) node[right=2pt,figlabel] {$G$};
  \fill[RadialTeal] (T) circle (.015) node[above left=1pt,figlabel] {$T$};
  \draw ($(T)+(-.025,-.008)$)--($(T)+(-.012,-.048)$)--($(T)+(.027,-.036)$);
  \node[figlabel,rotate=17] at (.05,.19) {$OT=\eta$};
  \node[figlabel,below] at (.62,-.03) {$OG=\rho$};
  \node[figlabel,rotate=-18] at (.69,.23)
    {$GT=\sqrt{\rho^2-\eta^2}$};
  \node[figlabel,align=center] at (.25,-1.48)
    {one side through $G$ is tangent to $\mathcal D_\eta$\\
     $\displaystyle\Lambda=
       \frac{3\eta+\sqrt{3(\rho^2-\eta^2)}}{2h}$};
\end{tikzpicture}%
}
\end{minipage}
\caption{The two minimizing configurations in the disk-plus-point formula.
For $\rho\le2\eta$, the centered equilateral triangle tangent to
$\mathcal D_\eta$ already contains $G$.  For $\rho>2\eta$, the point $G$
becomes a vertex of the minimizing triangle and one incident side is tangent
to the disk.  The tangent length
$GT=\sqrt{\rho^2-\eta^2}$ is the geometric source of the radical term.}
\label{fig:disk-plus-point-formula}
\end{figure}


\begin{proof}
This is the one-point specialization of
Proposition~\ref{prop:new-disk-finite-caliper}.  If $\rho\le2\eta$, orient
the three projections as $\rho/2,\rho/2,-\rho$.  Otherwise the minimum is a
point--disk tie.  The other active projection is
\[
 \frac{-\eta+\sqrt{3(\rho^2-\eta^2)}}2,
\]
which gives the displayed formula.
\end{proof}

For $p+q\le1$, put
\[
 J(p,q)=[X(q),X(1-p)]\subset e_{0,1},
 \qquad X(t)=V_0+t(V_1-V_0),
\]
\[
 c_*=c_{\max}(p,q),
 \qquad d=1-c_*.
\]

\begin{theorem}[Complementary-gap enclosure]
\label{thm:new-complementary-gap}
\[
 \boxed{
 \Lambda\left(\mathcal D_{hd}\cup J(p,q)\right)\ge1.
 }
 \tag{6.28}
\]
\end{theorem}

\begin{figure}[!htbp]
\centering
\begin{minipage}[t]{.315\linewidth}
\centering
\resizebox{\linewidth}{!}{%
\begin{tikzpicture}[scale=2.15,>=Latex]
  \node[panellabel,anchor=west] at (-1.05,1.16) {(a) the complementary gap};
  \HexCoords
  \DrawHexagon
  \draw[ray] (O)--($(V0)!.25!(V1)$);
  \coordinate (G) at ($(V0)!.25!(V1)$);
  \coordinate (H) at ($(V0)!.67!(V1)$);
  \draw[gap] (G)--(H);
  \fill[DemandRed] (G) circle (.025) node[right=2pt,figlabel] {$G=X(q)$};
  \fill[GapPurple] (H) circle (.021) node[left=2pt,figlabel] {$X(1-p)$};
  \draw[DemandRed,line width=.85pt] (O)--(G);
  \node[figlabel,align=center] at (0,-1.16)
    {$m=\min\{p,q\}$, choose the farther endpoint $G$\\
     $\lVert G\rVert^2=1-m+m^2$};
\end{tikzpicture}%
}
\end{minipage}\hfill
\begin{minipage}[t]{.335\linewidth}
\centering
\resizebox{\linewidth}{!}{%
\begin{tikzpicture}[scale=2.15,>=Latex]
  \node[panellabel,anchor=west] at (-1.08,1.16) {(b) six radial witnesses};
  \HexCoords
  \DrawHexagon
  \DrawRays
  \foreach \i in {0,...,5}{
    \coordinate (D\i) at ($(O)!.18!(V\i)$);
    \fill[DemandRed] (D\i) circle (.024);
  }
  \draw[DemandRed!65,line width=.75pt]
    (D0)--(D1)--(D2)--(D3)--(D4)--(D5)--cycle;
  \draw[CenterBlue,dashed,fill=CenterBlue!7,line width=.8pt]
    (O) circle ({.18*sqrt(3)/2});
  \node[figlabel,right] at (D0) {$D_0=dV_0$};
  \node[figlabel,above right] at (D1) {$D_1$};
  \node[figlabel,align=center] at (0,-1.16)
    {$d=1-c_*$, $c_*=c_{\max}(p,q)$\\
     $\mathcal D_{hd}\subset\operatorname{conv}\{D_0,\ldots,D_5\}$};
\end{tikzpicture}%
}
\end{minipage}\hfill
\begin{minipage}[t]{.315\linewidth}
\centering
\resizebox{\linewidth}{!}{%
\begin{tikzpicture}[x=2.25cm,y=2.25cm,>=Latex]
  \node[panellabel,anchor=west] at (-.72,1.02) {(c) enclosure terminal};
  \coordinate (O) at (0,0);
  \coordinate (G) at (1.02,.31);
  \draw[CenterBlue!35,fill=CenterBlue!8,line width=.8pt] (O) circle (.24);
  \fill[DemandRed] (G) circle (.023) node[right=2pt,figlabel] {$G$};
  \draw[black,line width=1pt]
    (-.42,-.64)--(1.18,-.18)--(.05,1.08)--cycle;
  \draw[DemandRed,dashed,line width=.75pt] (O)--(G);
  \node[figlabel] at (.03,-.08) {$\mathcal D_{hd}$};
  \node[figlabel,align=center] at (.23,-1.03)
    {$\displaystyle
      \Lambda\!\left(\mathcal D_{hd}\cup\{G\}\right)\ge1$\\
     hence $\Lambda(\{D_0,\ldots,D_5,G\})\ge1$};
\end{tikzpicture}%
}
\end{minipage}
\caption{The complementary-gap enclosure mechanism.  The farther endpoint
$G$ of $J(p,q)$ is center-forced.  Common-pair domination produces the six
center-forced points $D_i=(1-c_*)V_i$, whose convex hull contains the disk
$\mathcal D_{h(1-c_*)}$.  The disk-plus-point formula then forces enclosing
side length at least one.  The disk is not an additional witness; it is
contained in the convex hull of the six finite radial witnesses.}
\label{fig:complementary-gap-certificate}
\end{figure}


\begin{proof}
Let $m=\min\{p,q\}$.  The farther endpoint of the gap has squared norm
$\rho^2=1-m+m^2$.  Lemma~\ref{lem:new-disk-point-formula} reduces the claim
to
\[
 3c_*(1-c_*)\ge m(1-m).
 \tag{6.29}
\]
Common-pair domination gives $c_*\ge1-m\ge1/2$.  For $c_*\le h$, the left
side is at least $3h(1-h)>1/4$.  For $c_*>h$, the exact local support
inequalities give
\[
 m\le
 c_*\left(\frac12-\sqrt{c_*^2-\frac34}\right)=:\zeta(c_*).
\]
A direct calculation yields
\[
\begin{aligned}
&3c_*(1-c_*)-\zeta(c_*)(1-\zeta(c_*))\\
&\quad=
\frac{c_*(1-c_*)}{2}
\left(5-2c_*-2c_*^2+
2\sqrt{c_*^2-\frac34}\right)\ge0.
\end{aligned}
\]
Since $m,\zeta(c_*)<1/2$, this proves (6.29).
\end{proof}

\subsection{Universal terminal S: CE2 short ray}

\begin{theorem}[CE2 short-ray obstruction]
\label{thm:new-ce2-short-ray}
In the strict signed CE2 domain put
\[
 p=W-\alpha,\qquad q=R-\delta,\qquad
 e=\min\{\alpha,\delta\}.
\]
Then
\[
 \boxed{c_{\max}(p,q)<1-e.}
 \tag{6.30}
\]
Consequently, whenever both CE2 traces contain V-gaps and
$T_1,\ldots,T_5$ realize the common pair $(p,q)$, one of the radial
witnesses on $r_2,r_4$ lies beyond the corresponding C exit.
\end{theorem}

\begin{figure}[!htbp]
\centering
\begin{minipage}[t]{.31\linewidth}
\centering
\resizebox{\linewidth}{!}{%
\begin{tikzpicture}[scale=2.1,>=Latex]
  \node[panellabel,anchor=west] at (-1.08,1.18) {(a) two CE2 gaps};
  \HexCoords
  \DrawHexagon
  \DrawRays
  \draw[gap] ($(V0)!.18!(V1)$)--($(V0)!.36!(V1)$);
  \draw[gap] ($(V5)!.64!(V0)$)--($(V5)!.84!(V0)$);
  \draw[centerrole]
    (.07,.03)--(.45,.54)--(.62,-.35)--cycle;
  \node[figlabel,text=GapPurple] at (.55,.69) {demand $q$};
  \node[figlabel,text=GapPurple] at (.78,-.58) {demand $p$};
  \node[figlabel,align=center] at (0,-1.18)
    {$p,q$ are the common\\boundary demands};
\end{tikzpicture}%
}
\end{minipage}\hfill
\begin{minipage}[t]{.37\linewidth}
\centering
\resizebox{\linewidth}{!}{%
\begin{tikzpicture}[x=5cm,y=1cm,>=Latex]
  \node[panellabel,anchor=west] at (-.02,2.45) {(b) transverse center exits};
  \draw[flow] (0,1.55)--(1,1.55);
  \draw[flow] (0,.35)--(1,.35);
  \fill (0,1.55) circle (.018) node[left=2pt,figlabel] {$O$};
  \fill (1,1.55) circle (.018) node[right=2pt,figlabel] {$V_2$};
  \fill (0,.35) circle (.018) node[left=2pt,figlabel] {$O$};
  \fill (1,.35) circle (.018) node[right=2pt,figlabel] {$V_4$};
  \coordinate (E2) at (.08,1.55);
  \coordinate (D2) at (.14,1.55);
  \coordinate (D4) at (.14,.35);
  \coordinate (E4) at (.19,.35);
  \fill[CenterBlue] (E2) circle (.024) node[above=3pt,figlabel] {C trace end};
  \fill[DemandRed] (D2) circle (.024) node[below=3pt,figlabel] {$D_2$};
  \fill[DemandRed] (D4) circle (.024) node[above=3pt,figlabel] {$D_4$};
  \fill[CenterBlue] (E4) circle (.024) node[below=3pt,figlabel] {C trace end};
  \draw[GapPurple,line width=1.1pt] (E2)--(D2);
  \draw[black!35,line width=1.1pt] (D4)--(E4);
  \node[figlabel,align=center] at (.52,-.42)
    {one radial witness lies beyond the corresponding C trace};
\end{tikzpicture}%
}
\end{minipage}\hfill
\begin{minipage}[t]{.28\linewidth}
\centering
\resizebox{\linewidth}{!}{%
\begin{tikzpicture}[x=4.5cm,y=1cm,>=Latex]
  \node[panellabel,anchor=west] at (-.04,2.42) {(c) the case split};
  \node[figlabel,anchor=west] at (0,1.84) {first radial alternative};
  \draw[line width=.65pt] (0,1.38)--(1,1.38);
  \fill (0,1.38) circle (.017) node[left=1pt,figlabel] {$O$};
  \fill[CenterBlue] (.22,1.38) circle (.023) node[above=2pt,figlabel] {C trace};
  \fill[DemandRed] (.43,1.38) circle (.023) node[below=2pt,figlabel] {$D_2$};
  \draw[GapPurple,line width=1pt] (.22,1.38)--(.43,1.38);
  \node[figlabel,anchor=west] at (0,.73) {second radial alternative};
  \draw[line width=.65pt] (0,.27)--(1,.27);
  \fill (0,.27) circle (.017) node[left=1pt,figlabel] {$O$};
  \fill[CenterBlue] (.22,.27) circle (.023) node[above=2pt,figlabel] {C trace};
  \fill[DemandRed] (.43,.27) circle (.023) node[below=2pt,figlabel] {$D_4$};
  \draw[GapPurple,line width=1pt] (.22,.27)--(.43,.27);
  \node[figlabel,align=center] at (.5,-.36)
    {the red point lies beyond\\the corresponding center exit};
\end{tikzpicture}%
}
\end{minipage}
\caption{The CE2 short-ray terminal.  The two boundary gaps determine the
common pair $(p,q)$ and the center-forced radial witnesses $D_2,D_4$.  The
theorem proves that at least one of these witnesses lies beyond the
corresponding C-triangle trace and is uncovered.  The displayed ordering is
schematic.}
\label{fig:ce2-short-ray-certificate}
\end{figure}


\begin{proof}
Let $M=\max\{R,W\}$.  The CE2 inequalities
\[
 \alpha+W\delta<P,\qquad R\alpha+\delta<P
\]
imply
\[
 P>(1+M)e.
 \tag{6.31}
\]
Since $E^2=1-M+M^2$,
\[
 (1+M)^2-(3M-1)^2E^2
 =
 3M(1-M)(3M^2-2M+3)>0.
\]
Thus
\[
 \eta=P/E>(3M-1)e.
\]
Using $RW=\eta+P$ in the two CE2 inequalities gives
\[
 Wq>\eta+e>3Me,\qquad Rp>\eta+e>3Me.
\]
Because $R,W\le M$,
\[
 p,q>3e.
\]
Also
\[
 p+q=1-\alpha-\delta\le1-2e,
\]
so
\[
 \boxed{3e<p,q<1-5e.}
 \tag{6.32}
\]
Moreover
\[
 e<\frac{P}{1+M}
 \le\frac{2\sqrt3-3}{6}<\frac1{12}.
 \tag{6.33}
\]

Put $c=1-e$ and $f_c(t)=c^4-c^2+ct-t^2$.  The function is concave, and
\[
 f_c(3e)=e(1-7e-4e^2+e^3)>0,
\]
\[
 f_c(1-5e)=e(2-15e-4e^2+e^3)>0
\]
on the interval (6.33).  Therefore
\[
 f_c(p)>0,\qquad f_c(q)>0.
 \tag{6.34}
\]
Since $p+c,q+c>1$ and $c>p+q$, the four exact support-contact side lengths at
radial demand $c$ are
\[
 p+c,\quad q+c,\quad
 \frac{c^2}{\sqrt{p^2-pc+c^2}},\quad
 \frac{c^2}{\sqrt{q^2-qc+c^2}},
\]
and all four exceed one.  Hence $c$ is not admissible.

If $e=\delta$, the point $(1-c_{\max}(p,q))V_2$ lies beyond the C exit
$\delta$ on $r_2$.  If $e=\alpha$, the reflected conclusion holds on $r_4$.
\end{proof}

\UsesCloses{Lemmas~\ref{lem:readable-one-gap-common-pair},
\ref{lem:new-type-aware-radial-forcing}, and
\ref{lem:new-common-pair-domination}.}
{Every one-gap $N_+=0$ all-Vd0/T3-like row.}

\UsesCloses{Lemma~\ref{lem:readable-two-gap-common-pair} and the signed CE2
exit formula.}
{Every two-gap all-Vd0/T3-like row with $N_+\in\{0,1\}$.}
